% Entregable 1 — Consulta privada de inteligencia de amenazas mediante OT
% Pegar en Overleaf como proyecto nuevo. Compila con pdfLaTeX: pdfLaTeX -> BibTeX -> pdfLaTeX x2.
\documentclass[11pt,a4paper]{article}
\usepackage[a4paper,left=2cm,right=2cm,top=2.2cm,bottom=2.2cm]{geometry}
\usepackage[T1]{fontenc}
\usepackage[utf8]{inputenc}
\usepackage{lmodern}
\usepackage[spanish]{babel}
\AtBeginDocument{\shorthandoff{<>}}
\usepackage{amsmath,amssymb}
\usepackage{booktabs}
\usepackage{tabularx}
\usepackage{graphicx}
\usepackage{tikz}
\usetikzlibrary{arrows.meta,positioning,shapes}
\usepackage{algorithm}
\usepackage{algpseudocode}
\usepackage[hidelinks]{hyperref}
\usepackage[numbers,sort&compress]{natbib}

\setlength{\emergencystretch}{3em}
\renewcommand{\arraystretch}{1.15}
\newcommand{\code}[1]{\texttt{#1}}

\title{Consulta privada de inteligencia de amenazas mediante Oblivious Transfer\\ \large Entregable 1: Dise\~no y estado del arte (Proyecto 8)}
\author{Sebastian David Iba\~nez Rios \and Cesar Luis Roa Martinez \and Ivan Camilo Rueda Ariza \and Andres Arturo Carrero Martinez \\ \small Ingenieria de Sistemas --- Universidad del Norte}
\date{}

\begin{document}
\maketitle
\begin{abstract}
Un analista consulta 1 de $N$ indicadores de
amenazas sin revelar cu\'al. El proveedor publica un cat\'alogo versionado
de $N\in\{500,1000,2500,5000\}$ registros de 256\,B, cada uno cifrado con
clave independiente $K_j$ y AEAD. El cliente descarga el cat\'alogo cifrado
\emph{completo} (fase \emph{bulk}) para no filtrar su inter\'es y luego ejecuta
un OT 1-de-$N$ Diffie-Hellman (Boneh--Shoup \S11.6.1 + Remark~11.2) que entrega
solo la clave $K_i$. Privacidad del receptor perfecta, del emisor bajo CDH en
modelo de or\'aculo aleatorio, integridad v\'ia AEAD y HKDF con binding de
cat\'alogo, versi\'on, sesi\'on e \'indice. Se incluye l\'inea base HTTP y plan
de medici\'on del costo lineal $O(N)$ frente a la consulta convencional.
\end{abstract}

\tableofcontents
\newpage

\section{Escenario y contexto}
\label{sec:escenario}

\subsection{Caso de uso}
Un analista investiga una pista concreta (IoC: una IP, un hash o un n\'umero
de red ASN) del cat\'alogo de un proveedor porque apareci\'o en un equipo de
su empresa. Si la pide por HTTP con \code{GET /record/} y su identificador,
el proveedor ve exactamente qu\'e pista le interesa y puede deducir qu\'e
caso est\'a investigando. El objetivo es ocultar el \'indice $i$ elegido, y a
la vez garantizar que el cliente solo pueda abrir un registro por consulta.

El proveedor publica un cat\'alogo versionado de $N$ registros de longitud
fija (\textbf{256\,B} fijado; rango permitido 128--512) con
$N\in\{500,1000,2500,5000\}$. Cada registro se cifra con su propia clave
$K_j$ usando AEAD (cifrado que protege el contenido y detecta si alguien lo
modific\'o). El analista descarga el cat\'alogo cifrado \emph{completo} una
sola vez (fase \emph{bulk}, es decir, todo de una vez) y despu\'es abre
exactamente un registro con un protocolo OT 1-de-$N$ basado en
Diffie-Hellman, que solo le entrega la clave $K_i$ que necesita.

Los registros son sint\'eticos y de prueba: no son amenazas reales, solo
sirven para demostrar el protocolo. El Importador los crea con una semilla
fija (p.~ej. \code{seed=42}) para que cualquiera pueda repetir el mismo
cat\'alogo. Cada registro tiene seis campos con formato fijo: \code{record\_id}
como \code{ioc-2026-XXXX} seg\'un su posici\'on; \code{ip} rotando en el rango
de documentaci\'on \code{203.0.113.0/24}; \code{hash} como \code{sha256:} m\'as
64 d\'igitos hexadecimales aleatorios; \code{asn} como \code{AS64512} m\'as un
desplazamiento peque\~no dentro del rango privado de documentaci\'on;
\code{ts} como fecha ISO en UTC a partir de
\code{2026-09-25T11:00:00Z} sumando un minuto por registro; y \code{pad} como
texto de relleno.

Se serializa a JSON compacto con claves ordenadas y en UTF-8, y el largo de
\code{pad} se ajusta para que el JSON total mida exactamente 256\,B: no son
bytes pegados despu\'es, sino un campo de texto m\'as dentro del JSON (si el
contenido se pasa de 256\,B, el Importador da error y no recorta). As\'i todos
los registros quedan del mismo tama\~no y el tama\~no no filtra cu\'al se
pidi\'o. Despu\'es el Importador genera una clave $K_j$ por registro, cifra
ese JSON de 256\,B y guarda \code{catalog.json} + \code{keys.db}, como se
detalla en las Secciones~\ref{sec:arquitectura} y~\ref{sec:formatos}.

\begin{table}[ht]
\centering
\small
\begin{tabularx}{\textwidth}{@{}lX@{}}
\toprule
Campo & Formato y ejemplo \\
\midrule
\code{record\_id} & \code{ioc-2026-0042} (consecutivo seg\'un \code{idx}) \\
\code{ip} & \code{203.0.113.45} (rango de ejemplo, no real) \\
\code{hash} & \code{sha256:9f2c...} (64 hex aleatorios, aqu\'i recortado) \\
\code{asn} & \code{AS64512} (m\'as \code{idx mod 14}, rango privado) \\
\code{ts} & \code{2026-09-25T11:00:00Z} m\'as un minuto por \code{idx} \\
\code{pad} & texto de relleno (se alarga hasta que el JSON mida 256\,B) \\
\bottomrule
\end{tabularx}
\caption{Ejemplo de registro sint\'etico antes de cifrar. El JSON compacto ordenado, con \code{pad} incluido, mide siempre 256\,B.}
\end{table}

\subsection{Soluci\'on a alto nivel}
\begin{enumerate}
\item \textbf{Importador:} crea IoCs de prueba $\to$ los rellena hasta 256\,B
$\to$ genera una clave $K_j$ por registro y lo cifra
($C_j=\mathrm{AEAD.Enc}(K_j,\cdot)$) $\to$ guarda
\code{catalog.json} + \code{keys.db} (solo los tiene el servidor).
\item \textbf{Descarga bulk:} el cliente descarga $C=(C_1,\dots,C_N)$ entero.
Si pidiera un solo registro, el servidor sabr\'ia cu\'al le interesa.
\item \textbf{Sesi\'on OT:} en dos mensajes Diffie-Hellman el servidor le
entrega solo la clave $K_i$; el cliente descifra en su m\'aquina. Cada
consulta nueva usa un n\'umero de sesi\'on (\code{sid}) distinto.
\item \textbf{Baseline:} una consulta normal
\code{GET /record/} con identificador, para comparar tiempos y tama\~no de
datos, y as\'i medir cu\'anto cuesta ocultar la selecci\'on.
\end{enumerate}

\begin{center}
\begin{tikzpicture}[node distance=2.4cm, auto,
  box/.style={rectangle,draw,rounded corners,minimum width=2.6cm,minimum height=1cm,align=center}]
\node[box] (imp) {Importador};
\node[box,right=of imp] (srv) {Servidor\\(emisor OT)};
\node[box,right=of srv] (cli) {Cliente\\(CLI)};
\draw[-{Stealth}] (imp) -- node[above]{\scriptsize catalog.json} (srv);
\draw[-{Stealth}] (srv) -- node[above]{\scriptsize bulk $C$} (cli);
\draw[-{Stealth}] (cli) -- node[below]{\scriptsize OT$(i)\to K_i$} (srv);
\end{tikzpicture}
\end{center}

\section{Modelo de adversario}
\label{sec:adversario}
Aqu\'i definimos contra qui\'en nos defendemos: qu\'e puede intentar cada
parte y qu\'e queda fuera del modelo.

\begin{table}[ht]
\centering
\small
\begin{tabularx}{\textwidth}{@{}l l X@{}}
\toprule
Actor & Se porta as\'i & Qu\'e intenta \\
\midrule
Servidor & Cumple el protocolo pero es curioso &
Sigue los pasos, pero intenta adivinar el $i$ mirando los mensajes
$(u,\{E_j\})$. No lo logra si el problema DH es dif\'icil. \\
Cliente & Quiere m\'as de lo permitido &
Cumple para pedir, pero intenta sacar $K_j$ de\'indices que no pidi\'o
($j\ne i$) con una sola respuesta. No puede sin romper el problema DH. \\
Red & Esp\'ia y modifica &
Mira el tr\'afico o cambia bytes. Cualquier cambio debe ser detectado
por el cifrado AEAD y el control de versi\'on. \\
\bottomrule
\end{tabularx}
\caption{Modelo de adversario del sistema.}
\end{table}

\textbf{Supuestos y l\'imites del modelo}:
(i)~PIR, extensi\'on de OT, OT adaptativo con OPRF ni bases de millones de
registros;
(ii)~clientes o servidores totalmente tramposos (solo cubrimos servidor
curioso que cumple el protocolo y cliente que pide de m\'as; si el servidor
manda datos falsos solo lo detectamos por integridad, como en el ejemplo de
los art\'iculos vac\'ios de la Secci\'on~11.6.4 del libro);
(iii)~pruebas de seguridad para cualquier combinaci\'on de programas
(componibilidad);
(iv)~ocultar IP, frecuencia, momento o volumen de consultas, ni pagos o
licenciamiento real.

\section{Propiedades de seguridad}
\label{sec:propiedades}
Nombres que usamos desde aqu\'i (ver Secci\'on~\ref{sec:protocolo}):
$N$ es el tama\~no del cat\'alogo, $i\in[1,N]$ el \'indice elegido,
$K_j$ la clave del registro $j$,
$\mathsf{ctx}$ el contexto inmutable que concatena de forma can\'onica
$\mathit{catalog\_id}$, $\mathit{version}$ y $\mathit{group\_id}$, y
$sid\in\{0,1\}^{128}$ el identificador \'unico de sesi\'on fresca.

El sistema debe satisfacer formalmente tres propiedades:

\begin{enumerate}
\item \textbf{Privacidad del receptor (selecci\'on oculta perfecta):}
Para cualesquiera dos \'indices $i_0, i_1 \in \{1,\dots,N\}$, la vista del
servidor es perfectamente indistinguible:
$\mathrm{Vista}_S(i_0) \equiv \mathrm{Vista}_S(i_1)$.
En nuestro protocolo, el mensaje $u = g^\alpha \cdot v^{-i}$ con $\alpha$
escogido uniformemente al azar en $\mathbb{Z}_q$ est\'a uniformemente
distribuido sobre el grupo $G$, enmascarando $i$ como una libreta de un solo
uso en el exponente. El servidor aprende exactamente cero informaci\'on sobre
$i$ en el sentido de la teor\'ia de la informaci\'on.

\item \textbf{Privacidad del emisor (transferencia \'unica bajo CDH):}
El cliente aprende $K_i$ y no obtiene informaci\'on computable sobre ninguna de
las dem\'as claves $\{K_{j\ne i}\}$.
Si un cliente malicioso pudiera extraer dos claves distintas en una sola
ejecuci\'on, podr\'ia calcular $v^\beta = g^{\beta^2}$, lo cual es
computacionalmente equivalente a resolver el problema Computational
Diffie-Hellman (CDH) en $G$ (seg\'un el Ejercicio~10.17 del libro
\cite{boneh-shoup-v06}), adem\'as de tener que romper el or\'aculo aleatorio
asociado a HKDF.

\item \textbf{Integridad, autenticidad y vinculaci\'on de contexto:}
Cualquier modificaci\'on no autorizada es detectada inmediatamente por el
cifrado autenticado (AEAD), arrojando la excepci\'on \code{InvalidTag}. Para
evitar ataques de confusi\'on de dominio o de traslaci\'on entre sesiones, se
definen dos datos asociados ($AAD$) diferenciados:
\begin{itemize}
\item \textbf{Dato asociado del cat\'alogo:}
$AAD_{\mathrm{cat}, j} = \mathsf{ctx}\Vert \mathrm{format\_u32}(j)$, utilizado
al cifrar cada registro $C_j$ de forma offline (no contiene $sid$ porque el
cat\'alogo es est\'atico e independiente de las sesiones).
\item \textbf{Dato asociado de la sesi\'on OT:}
$AAD_{\mathrm{ot}, j} = \mathsf{ctx}\Vert sid\Vert \mathrm{format\_u32}(j)$,
utilizado al encapsular cada clave $E_j$, ligando la respuesta a la sesi\'on
activa y al \'indice de posici\'on.
\end{itemize}
Asimismo, la huella $\mathrm{sha256}$ de la colecci\'on de textos cifrados
permite al cliente verificar que el cat\'alogo no haya sido alterado antes de
iniciar cualquier consulta.
\end{enumerate}
Las justificaciones matem\'aticas detalladas se presentan en la
Secci\'on~\ref{sec:protocolo}.

\section{Estado del arte y protocolo derivado}
\label{sec:protocolo}

\subsection{Marco te\'orico y an\'alisis del estado del arte}
El protocolo de transferencia inconsciente 1-de-$N$ (\emph{1-out-of-$N$ OT})
permite que un emisor con $N$ secretos $m_1,\dots,m_N$ interact\'ue con un
receptor que posee un \'indice privado $i \in \{1,\dots,N\}$, de modo que al
finalizar la ejecuci\'on el receptor aprenda $m_i$ sin obtener informaci\'on
alguna sobre los restantes $N-1$ mensajes, mientras que el emisor no aprende
absolutamente nada sobre el valor de $i$ \cite{boneh-shoup-v06}.

En la literatura existen m\'ultiples aproximaciones para construir esquemas OT:
\begin{itemize}
\item \textbf{Construcci\'on pedag\'ogica de Boneh y Shoup (\S11.6.1)
\cite{boneh-shoup-v06}:} Propone un OT 1-de-$N$ derivado del cifrado ElGamal
en un grupo c\'iclico de orden primo $G=\langle g\rangle$. El emisor sortea
$\beta$ y publica $v = g^\beta$; el receptor camufla su \'indice calculando
$u = g^\alpha \cdot v^{-i}$ con un escalar fresco $\alpha$; el emisor deriva
$N$ claves p\'ublicas ElGamal impl\'icitas $u_j = u \cdot v^j$ y encripta cada
mensaje reutilizando $\beta$. Con las salvaguardas de sesi\'on y AEAD sugeridas
en el Remark~11.2, esta construcci\'on es limpia, rigurosa y verificable,
siendo la referencia obligatoria adoptada en este proyecto.
\item \textbf{Naor y Pinkas (2001, 2005) \cite{naor-pinkas-soda01,naor-pinkas-joc05}:}
Permiten construir un OT 1-de-$N$ a partir de $\lceil \log_2 N \rceil$
instancias paralelas de OT 1-de-2. Aunque optimizan el ancho de banda asint\'otico,
introducen m\'as rondas de negociaci\'on y mayor complejidad de estado para
el rango moderado de nuestro cat\'alogo ($N \in [500, 5000]$).
\item \textbf{Tzeng (2002) \cite{tzeng-pkc02}:} Propone un esquema directo de
una sola ronda; sin embargo, para garantizar la seguridad frente a receptores
maliciosos requiere pruebas de conocimiento cero (ZKP) sobre la construcci\'on
del mensaje del cliente, elevando la sobrecarga algor\'itmica.
\item \textbf{Chou y Orlandi (2015) \cite{chou-orlandi-simplestOT}:} Dise\~nan
un OT de alta velocidad (\emph{The Simplest OT}), pero su implementaci\'on en
curvas el\'ipticas con cofactor $h>1$ exige un manejo cuidadoso de subgrupos
para evitar fugas de bits, lo que motiva el uso de \texttt{ristretto255}.
\item \textbf{OT adaptativo y OPRF de Boneh--Shoup (\S11.6.3--\S11.6.4):}
Permiten evaluar funciones pseudoaleatorias inconscientes con costo de
comunicaci\'on $O(1)$ por consulta; no obstante, el enunciado del proyecto las
declara expresamente \emph{fuera de alcance}, exigiendo el protocolo Diffie-Hellman
lineal pedag\'ogico.
\end{itemize}

Por estas razones, se descartan las alternativas anteriores y se adopta la
construcci\'on de Boneh--Shoup \S11.6.1 con las extensiones del Remark~11.2.

\subsection{Especificaci\'on del protocolo adaptado}
Sean $G = \langle g \rangle$ un grupo de orden primo $q$ (\texttt{ristretto255}),
$\mathrm{HKDF}$ la funci\'on de derivaci\'on HKDF-SHA256 \cite{rfc5869-hkdf} y
$\mathrm{AEAD}$ el cifrado AES-256-GCM \cite{nist-gcm}. Se define la cadena de
contexto global $\mathsf{ctx}$ (ver Secci\'on~\ref{sec:formatos}).

El protocolo se ejecuta en las siguientes fases:

\begin{algorithm}[ht]
\caption{Protocolo interactivo OT 1-de-$N$ para entrega privada de $K_i$}
\label{alg:ot-protocol}
\small
\begin{algorithmic}[1]
\Require \textbf{Servidor:} claves $\{K_j\}_{j=1}^N$, contexto $\mathsf{ctx}$. \textbf{Cliente:} \'indice $i \in \{1,\dots,N\}$, $\{C_j\}_{j=1}^N$, $\mathsf{ctx}$.
\Ensure \textbf{Cliente:} registro descifrado $record_i$; el servidor no aprende $i$.
\State \textbf{Fase 0 (Offline):} $C_j \gets \mathrm{AEAD.Enc}(K_j, \mathit{nonce}_j, record_j, \mathsf{ctx}\Vert \mathrm{format\_u32}(j))$. El cliente descarga $C$ entero.
\State \textbf{Paso 1 (Inicio de sesi\'on):} Cliente invoca \code{POST /ot/session}. Servidor sortea $\beta \xleftarrow{R} \mathbb{Z}_q$, calcula $v \gets g^\beta \in G$,
\State \quad genera $sid \xleftarrow{R} \{0,1\}^{128}$, guarda $(sid, \beta)$ (TTL 60\,s) y responde $(sid, v)$.
\State \textbf{Paso 2 (Consulta OT):} Cliente sortea $\alpha \xleftarrow{R} \mathbb{Z}_q$, calcula $u \gets g^\alpha \cdot v^{-i} \in G$, env\'ia $(sid, u)$ a \code{POST /ot/query}.
\State \textbf{Paso 3 (Respuesta del emisor):} Servidor extrae y destruye $\beta$ para $sid$. Valida $u \in G$. Para $j = 1 \dots N$:
\State \quad $u_j \gets u \cdot v^j$ ($u_0 = u, u_j = u_{j-1}\cdot v$); $Z_j \gets (u_j)^\beta$;
\State \quad $k_j \gets \mathrm{HKDF}(\mathsf{ctx}, Z_j, \texttt{oblivion-ti/ot-v1}\Vert sid \Vert \mathrm{format\_u32}(j))$;
\State \quad $E_j \gets \mathrm{AEAD.Enc}(k_j, \mathit{nonce}'_j, K_j, \mathsf{ctx}\Vert sid \Vert \mathrm{format\_u32}(j))$. Responde con $\{(\mathit{nonce}'_j, E_j)\}_{j=1}^N$.
\State \textbf{Paso 4 (Descifrado local):} Cliente calcula $Z_i \gets v^\alpha = g^{\alpha\beta}$; deriva $k_i$;
\State \quad abre $K_i \gets \mathrm{AEAD.Dec}(k_i, \mathit{nonce}'_i, E_i, \mathsf{ctx}\Vert sid \Vert \mathrm{format\_u32}(i))$;
\State \quad abre $record_i \gets \mathrm{AEAD.Dec}(K_i, \mathit{nonce}_i, C_i, \mathsf{ctx}\Vert \mathrm{format\_u32}(i))$.
\end{algorithmic}
\end{algorithm}

En t\'erminos de red sobre HTTP, la interacci\'on comprende:
(1) descarga en bloque de \code{catalog.json};
(2) solicitud de sesi\'on (\code{POST /ot/session} $\to (sid, v)$); y
(3) consulta OT (\code{POST /ot/query} $\to (sid, u)$ y respuesta con $\{E_j\}$).

\begin{figure}[ht]
\centering
\small
\begin{tikzpicture}[node distance=0.8cm, auto,
  party/.style={rectangle,draw,minimum width=2.2cm,minimum height=0.7cm,align=center,font=\footnotesize},
  msg/.style={font=\scriptsize}]
\node[party] (r) {Receptor\\(Cliente $i$)};
\node[party,right=4.2cm of r] (s) {Emisor\\(Servidor)};
\draw[thick] (r.south) -- ++(0,-3.2);
\draw[thick] (s.south) -- ++(0,-3.2);

\draw[-{Stealth}] ([yshift=-0.5cm]r.south) -- node[above,msg] {\code{POST /ot/session}} ([yshift=-0.5cm]s.south);
\draw[-{Stealth}] ([yshift=-1.1cm]s.south) -- node[above,msg] {$(sid, v = g^\beta)$} ([yshift=-1.1cm]r.south);
\draw[-{Stealth}] ([yshift=-1.8cm]r.south) -- node[above,msg] {\code{POST /ot/query}: $(sid, u = g^\alpha v^{-i})$} ([yshift=-1.8cm]s.south);
\draw[-{Stealth}] ([yshift=-2.6cm]s.south) -- node[above,msg] {$\{E_j, \mathit{nonce}'_j\}_{j=1}^N$} ([yshift=-2.6cm]r.south);
\end{tikzpicture}
\caption{Diagrama de secuencia e interacci\'on de la sesi\'on OT.}
\label{fig:diagrama-secuencia-ot}
\end{figure}

\subsection{Manejo de sesiones repetidas e interrumpidas}
El dise\~no resuelve de forma integral las sesiones interrumpidas y repetidas:
\begin{itemize}
\item \textbf{Sesiones interrumpidas:} Cuando el servidor genera $(sid, \beta, v)$
en el Paso~1, almacena el par en una tabla en memoria con un tiempo l\'imite de
vida (\emph{TTL}) fijado en 60 segundos. Si el cliente interrumpe la ejecuci\'on,
se desconecta o sufre un fallo de red antes de enviar el Paso~2, la entrada
caduca y $\beta$ es purgado autom\'aticamente, impidiendo fugas de memoria o
acumulaci\'on de estados obsoletos.
\item \textbf{Sesiones de un solo uso (\emph{single-use}):} Al procesar la consulta
en el Paso~3, el servidor extrae $\beta$ y lo destruye de inmediato. Si un
adversario o cliente reintenta enviar el mismo $sid$ (\emph{replay attack}),
el servidor rechaza la petici\'on con error HTTP \code{409 Conflict} o
\code{404 Not Found}.
\end{itemize}

\subsection{Justificaci\'on rigurosa de seguridad}

\subsubsection{Privacidad del receptor (ocultamiento perfecto de $i$)}
La \'unica informaci\'on que el servidor recibe del cliente dependiente del
\'indice $i$ es el elemento de grupo $u = g^\alpha \cdot v^{-i} \in G$.
Dado que $G$ es un grupo c\'iclico de orden primo $q$ y que $\alpha$ es elegido
por el cliente de forma estrictamente uniforme en $\mathbb{Z}_q$, la aplicaci\'on
$\alpha \mapsto g^\alpha$ es un isomorfismo que distribuye $g^\alpha$ de forma
uniforme sobre $G$. Como $v^{-i}$ es un elemento fijo de $G$, la operaci\'on
$u = g^\alpha \cdot v^{-i}$ constituye una traslaci\'on biyectiva sobre $G$,
actuando como una libreta de un solo uso (\emph{one-time pad}) en el exponente.
Por lo tanto, para cualquier valor fijado de $v$ e $i$, $u$ es una variable
aleatoria uniformemente distribuida sobre $G$. La vista del servidor es
estad\'isticamente independiente de $i$, garantizando privacidad incondicional
(te\'orica de la informaci\'on).

\subsubsection{Privacidad del emisor (reducci\'on al problema CDH)}
Para demostrar que el cliente no puede aprender m\'as de una clave sim\'etrica
por ejecuci\'on, consideremos un adversario $\mathcal{A}$ capaz de derivar dos
claves distintas $K_{j_1}$ y $K_{j_2}$ con $j_1 \ne j_2$. Bajo el modelo de
or\'aculo aleatorio para $\mathrm{HKDF}$, $\mathcal{A}$ debe haber consultado
los dos secretos $Z_{j_1}$ y $Z_{j_2}$, los cuales satisfacen:
\begin{equation}
Z_{j_1} = (u \cdot v^{j_1})^\beta = u^\beta \cdot (v^\beta)^{j_1}, \qquad
Z_{j_2} = (u \cdot v^{j_2})^\beta = u^\beta \cdot (v^\beta)^{j_2}
\end{equation}
Dividiendo ambos elementos en el grupo $G$:
\begin{equation}
\frac{Z_{j_2}}{Z_{j_1}} = \frac{u^\beta \cdot (v^\beta)^{j_2}}{u^\beta \cdot (v^\beta)^{j_1}} = (v^\beta)^{j_2 - j_1}
\end{equation}
Puesto que $j_1 \ne j_2$ y $1 \le |j_2 - j_1| < N < q$, la diferencia $(j_2 - j_1)$
es invertible en $\mathbb{Z}_q$. Calculando la potencia con el inverso modular
$\Delta = (j_2 - j_1)^{-1} \pmod q$:
\begin{equation}
\left( \frac{Z_{j_2}}{Z_{j_1}} \right)^\Delta = \left( (v^\beta)^{j_2 - j_1} \right)^{(j_2 - j_1)^{-1}} = v^\beta = (g^\beta)^\beta = g^{\beta^2}
\end{equation}
Por el Teorema~10.1 y el Ejercicio~10.17 de Boneh y Shoup \cite{boneh-shoup-v06},
el problema de calcular $g^{\beta^2}$ a partir de $g^\beta$ (\emph{Squaring Diffie-Hellman})
es computacionalmente equivalente al problema Computational Diffie-Hellman (CDH).
En consecuencia, la existencia de un receptor capaz de abrir dos claves en una
misma sesi\'on permitir\'ia romper el supuesto CDH en $G$.

\subsubsection{Mitigaci\'on del ataque entre sesiones (Ejercicio 11.20 y Remark 11.2)}
Si el servidor reusara el mismo par $(\beta, v)$ en m\'ultiples consultas sin
vinculaci\'on de sesi\'on, un cliente que conoci\'o $Z_i = v^\alpha$ para la
consulta de $i$ podr\'ia enviar en una segunda consulta el valor manipulado
$\hat{u} = u \cdot v$. El servidor computar\'ia $\hat{u}_j = \hat{u}\cdot v^j = u\cdot v^{j+1} = u_{j+1}$,
lo que provocar\'ia que $\hat{Z}_{i-1} = Z_i$. El cliente podr\'ia descifrar
la encapsulaci\'on $\hat{E}_{i-1}$ utilizando el secreto que ya conoc\'ia,
aprendiendo un segundo registro sin autorizaci\'on.

Esta vulnerabilidad queda totalmente neutralizada en nuestro dise\~no:
\begin{enumerate}
\item \textbf{Escalar ef\'imero fresco:} Cada sesi\'on $sid$ genera un $\beta$
nuevo e independiente, haciendo que $v$ sea impredecible y no correlacionable.
\item \textbf{Separaci\'on de dominios en HKDF:} La derivaci\'on incluye
$sid$ e \'indice $j$ en el par\'ametro \code{info}, produciendo claves $k_j$
disjuntas aun si hubiera coincidencias algebraicas en $Z_j$.
\item \textbf{Autenticaci\'on v\'ia AEAD:} Las encapsulaciones $E_j$ incluyen
$AAD_{\mathrm{ot}, j} = \mathsf{ctx}\Vert sid\Vert \mathrm{format\_u32}(j)$.
Cualquier intento de abrir un $E_j$ de otra sesi\'on falla deterministamente
con \code{InvalidTag}.
\end{enumerate}

\subsection{Conteo exacto de operaciones y costos}
El conteo preciso de operaciones algebraicas en el grupo $G$ es:
\begin{itemize}
\item \textbf{Cliente (Receptor):} Realiza exactamente \textbf{3 exponenciaciones}
(c\'alculo de $g^\alpha$, c\'alculo de $v^{-i}$ y c\'alculo de $Z_i = v^\alpha$)
y \textbf{1 multiplicaci\'on} ($g^\alpha \cdot v^{-i}$). Su costo computacional
es $O(1)$.
\item \textbf{Servidor (Emisor):} Realiza exactamente \textbf{$N+1$ exponenciaciones}
(1 para $v = g^\beta$ y $N$ para calcular cada $Z_j = u_j^\beta$) y \textbf{$N$ multiplicaciones}
de grupo (reconstrucci\'on de $u_j = u_{j-1}\cdot v$ con $u_0 = u$). Su costo
computacional es $O(N)$.
\item \textbf{Sobrecarga de red:} El cliente transmite 48\,B en la petici\'on
($sid$ de 16\,B y $u$ de 32\,B). El servidor transmite $48 + 60N$\,bytes en la
respuesta ($sid$, $v$ y $N$ encapsulaciones de 48\,B con nonces de 12\,B).
\end{itemize}

\section{Arquitectura del sistema}
\label{sec:arquitectura}

\textbf{Importador} (programa offline en Python): crea IoCs de prueba con
relleno hasta 256\,B, genera las claves $K_j$ y los cifrados $C_j$, y guarda
\code{catalog.json} m\'as \code{keys.db} (esto \'ultimo solo lo tiene el
servidor).
\textbf{Servidor} (el que entrega, con \code{FastAPI} y \code{uvicorn}):
entrega la descarga completa y responde las peticiones OT, adem\'as del
endpoint de comparaci\'on.
\textbf{Cliente} (programa de consola con \code{Typer}):
descarga, pide por OT, descifra y modo de comparaci\'on.
Flujo general: descarga completa una vez y luego una sesi\'on OT nueva por
cada registro que se quiera abrir.

\begin{table}[ht]
\centering
\small
\begin{tabularx}{\textwidth}{@{}l l X@{}}
\toprule
Ruta & M\'etodo & Para qu\'e sirve \\
\midrule
\code{/catalog} & \code{GET} & Descarga completa de \code{catalog.json}. \\
\code{/ot/query} & \code{POST} & Se env\'ia $(u,sid)$ y devuelve $E_{1..N}$. \\
\code{/record} con id & \code{GET} & Solo comparaci\'on normal (baseline). \\
\bottomrule
\end{tabularx}
\caption{Rutas HTTP m\'inimas. TODO(equipo): diagrama de despliegue local.}
\end{table}

\textbf{Comandos del cliente:}
\code{fetch-catalog} (descargar), \code{ot-request} (pedir por OT),
\code{decrypt} (abrir) y \code{baseline-get} (pedido normal).

\textbf{L\'inea base:} la misma aplicaci\'on pero pidiendo el registro
directo por HTTP, sin OT. Sirve para comparar tiempo total, datos
transferidos y operaciones, y as\'i medir cu\'anto cuesta ocultar $i$.

\section{Formatos y serializaci\'on}
\label{sec:formatos}

El cat\'alogo y los mensajes son JSON. Todo valor binario (nonces, textos
cifrados y elementos de grupo) se codifica en \code{base64url} sin relleno
(RFC~4648, \S5 \cite{rfc4648}). Los elementos de \texttt{ristretto255} viajan
en su codificaci\'on can\'onica de 32\,B \cite{rfc9496}.

\begin{table}[ht]
\centering
\small
\begin{tabularx}{\textwidth}{@{}l X@{}}
\toprule
Objeto & Campos y tipos exactos \\
\midrule
\code{catalog.json} &
\code{catalog\_id} (texto), \code{version} (texto, p.~ej. ``v1''),
\code{group\_id} (texto, ``ristretto255''), \code{N} (entero),
\code{record\_len} $=256$, \code{sha256\_root} (hex de 64 caracteres) y
\code{records}: lista ordenada de objetos con \code{record\_id} (texto),
\code{idx} ($1..N$), \code{nonce} (b64url de 12\,B) y \code{ct} (b64url de 272\,B:
registro cifrado m\'as tag). \\
\code{SessionRequest} &
\code{catalog\_version} (texto, p.~ej. ``v1''). \\
\code{SessionResponse} &
\code{sid} (hex de 32 caracteres, 16\,B), \code{v\_b64} (punto $v$ en b64url de 32\,B) y
\code{ttl} (entero, 60 segundos). \\
\code{OTRequest} &
\code{sid} (hex de 16\,B) y \code{u\_b64} (punto $u = g^\alpha v^{-i}$ en b64url de 32\,B).
\textbf{El \'indice $i$ nunca viaja en claro ni impl\'icito.} \\
\code{OTResponse} &
\code{sid} (hex) y \code{encaps}: lista ordenada de $N$ objetos con \code{idx}
($1..N$), \code{nonce} (b64url de 12\,B) y \code{ct} (b64url de 48\,B, clave $K_j$
cifrada con tag). \\
\bottomrule
\end{tabularx}
\caption{Formatos de persistencia y mensajes de red.}
\label{tab:formatos-mensajes}
\end{table}

\subsection{Codificaci\'on can\'onica y huella}
\begin{enumerate}
\item \textbf{Contexto $\mathsf{ctx}$:} cada cadena va precedida de su
longitud en 2 bytes sin signo (\emph{big-endian}), para que la
concatenaci\'on no sea ambigua:
\begin{equation}
\mathsf{ctx} = \mathrm{u16}(|\mathit{cat\_id}|)\Vert \mathit{cat\_id} \Vert
\mathrm{u16}(|\mathit{ver}|)\Vert \mathit{ver} \Vert
\mathrm{u16}(|\mathit{group}|)\Vert \mathit{group}
\end{equation}
\item \textbf{$sid$ e \'indice $j$:} $sid$ se usa como sus 16 bytes
(\code{bytes.fromhex(sid)}) y el \'indice como entero sin signo de 32 bits
en orden de red: $\mathrm{format\_u32}(j) = j.\mathrm{to\_bytes}(4, \text{'big'})$.
Al ser ambos de longitud fija, las concatenaciones siguientes tampoco son
ambiguas.
\item \textbf{Datos asociados:}
$AAD_{\mathrm{cat}, j} = \mathsf{ctx} \Vert \mathrm{format\_u32}(j)$ protege
el registro en reposo, y
$AAD_{\mathrm{ot}, j} = \mathsf{ctx} \Vert sid \Vert \mathrm{format\_u32}(j)$
liga la clave transferida a la sesi\'on y a la posici\'on $j$.
\item \textbf{Huella \code{sha256\_root}:} SHA-256 de la concatenaci\'on, en
orden ascendente de $j$, de los bytes crudos de cada registro del cat\'alogo:
\begin{equation}
\code{sha256\_root} = \mathrm{SHA\text{-}256}\big(\mathit{nonce}_1 \Vert C_1
\Vert \mathit{nonce}_2 \Vert C_2 \Vert \cdots \Vert \mathit{nonce}_N \Vert C_N\big)
\end{equation}
El cliente la recalcula tras descargar \code{catalog.json}; si no coincide,
aborta.
\end{enumerate}

\subsection{Reglas de validaci\'on y rechazo}
\begin{itemize}
\item \textbf{Elementos de grupo mal formados:} si los 32 bytes de $u$ (en
el servidor) o de $v$ (en el cliente) no son una codificaci\'on can\'onica
v\'alida de \texttt{ristretto255} (reglas de decodificaci\'on de RFC~9496,
\S4.3.1; vectores de prueba en su Ap\'endice~A.2), se rechazan: el servidor
responde \code{400 Bad Request} y el cliente aborta. Lo mismo ocurre si $v$ o
alg\'un $u_j$ es el elemento identidad, que la biblioteca se niega a
multiplicar.
\item \textbf{\'Indices inv\'alidos:} el \'indice nunca llega al servidor,
as\'i que se valida solo en el cliente: si el \code{record\_id} no est\'a en
el mapa p\'ublico o $i \notin [1,N]$, el cliente aborta antes de abrir
sesi\'on. Tambi\'en rechaza una \code{OTResponse} que no traiga exactamente
$N$ entradas con \code{idx} $= 1,\dots,N$ en orden.
\item \textbf{Sesiones duplicadas o expiradas:} un $sid$ ya consumido recibe
\code{409 Conflict}; uno desconocido o vencido, \code{404 Not Found}.
\item \textbf{Versi\'on incompatible:} si \code{catalog\_version} no coincide
con la del servidor, este responde \code{422 Unprocessable Entity} y el
cliente debe volver a descargar el cat\'alogo.
\item \textbf{Mapa \code{record\_id}--\code{idx}:} es p\'ublico, va dentro de
\code{catalog.json} y no cambia dentro de una misma versi\'on.
\end{itemize}

\section{Selecci\'on de bibliotecas y primitivas}
\label{sec:primitivas}

\begin{table}[ht]
\centering
\small
\begin{tabularx}{\textwidth}{@{}l l X@{}}
\toprule
Para qu\'e & Qu\'e usamos & Por qu\'e \\
\midrule
Grupo matem\'atico &
\textbf{ristretto255} (RFC~9496 \cite{rfc9496}) &
Grupo c\'iclico de orden primo $q \approx 2^{252}$. Elimina el cofactor 8 de
Curve25519, con lo que desaparecen los ataques de subgrupo peque\~no
\cite{chou-orlandi-simplestOT}. Cada elemento tiene una \'unica codificaci\'on
can\'onica de 32\,B y la decodificaci\'on rechaza las dem\'as. \\
Derivar claves &
\textbf{HKDF-SHA256} (RFC~5869 \cite{rfc5869-hkdf}) &
Convierte el secreto Diffie-Hellman $Z_j$ en una clave de 256 bits. La sal
$\mathsf{ctx}$ separa cat\'alogos y versiones; el \code{info} ata el dominio
\code{"oblivion-ti/ot-v1"}, el $sid$ y el \'indice $j$. \\
Cifrado sim\'etrico &
\textbf{AES-256-GCM} (NIST SP 800-38D \cite{nist-gcm}) &
Cifrado autenticado (AEAD) con tag de 128 bits y nonce aleatorio de 96 bits.
Cada clave ($K_j$ y $k_j$) cifra un solo mensaje, as\'i que no hay riesgo de
repetir nonce. El dato asociado ata contexto, sesi\'on e \'indice. \\
\bottomrule
\end{tabularx}
\caption{Primitivas criptogr\'aficas seleccionadas y su justificaci\'on.}
\label{tab:seleccion-primitivas}
\end{table}

\subsection{Bibliotecas}
AES-256-GCM y HKDF-SHA256 se toman de \code{cryptography} (sobre OpenSSL,
con aceleraci\'on por hardware). Para el grupo, el protocolo necesita tres
operaciones sobre elementos arbitrarios: producto ($u_j = u_{j-1}\cdot v$),
exponenciaci\'on ($u_j^\beta$, $v^{-i}$) y validaci\'on de codificaciones.
\code{cryptography} no las ofrece: su API de curvas el\'ipticas se limita a
ECDH y ECDSA y no permite sumar dos puntos. \code{PyNaCl} tampoco: se
comprob\'o en las versiones 1.5.0 y 1.6.2 que expone esa aritm\'etica para
Ed25519, cuyo grupo tiene cofactor 8, pero no las funciones de
\texttt{ristretto255}. Se adopta por eso \code{rbcl}, que empaqueta
\code{libsodium} y expone sus funciones \code{crypto\_core\_ristretto255\_*}
(producto, validaci\'on de puntos y aritm\'etica de escalares) y
\code{crypto\_scalarmult\_ristretto255} (exponenciaci\'on). Como el producto
no valida sus entradas, todo elemento recibido se comprueba antes con
\code{is\_valid\_point}. Como respaldo queda NIST P-256, tambi\'en de orden
primo, con \code{PyCryptodome}, cuya clase \code{EccPoint} permite sumar y
multiplicar puntos; solo cambiar\'ian \code{group\_id} y la serializaci\'on
de los elementos.

Las dependencias se fijan a versiones exactas en \code{requirements.txt},
sobre Python~3.11: rbcl 1.1.2, cryptography 50.0.2, fastapi 0.142.2, uvicorn
0.54.0, typer 0.27.2, pydantic 2.13.5 y httpx 0.28.1 (cliente HTTP).

\section{Plan de pruebas y experimentaci\'on}
\label{sec:pruebas}

Este plan define c\'omo se demostrar\'a la correcci\'on funcional, la
seguridad frente a manipulaciones y el desempe\~no del sistema; la
implementaci\'on y las mediciones se presentar\'an en el Entregable~2. Todas
las pruebas del Cuadro~\ref{tab:plan-pruebas} ser\'an automatizadas.

\begin{table}[ht]
\centering
\small
\begin{tabularx}{\textwidth}{@{}l p{6.2cm} X@{}}
\toprule
Tipo & Acci\'on de prueba & Comportamiento esperado \\
\midrule
Casos v\'alidos &
Pedir el primer registro ($i=1$), el \'ultimo ($i=N$) y varios al azar;
repetir el mismo $i$ con una sesi\'on nueva. &
El registro descifrado coincide byte a byte con el original; $sid$, $u$, $v$
y $\{E_j\}$ son distintos en cada ejecuci\'on. \\
Claves no pedidas &
Intentar abrir $E_{j\ne i}$ con la clave derivada de $Z_i$. &
AEAD rechaza con \code{InvalidTag}. \\
Entradas inv\'alidas &
Enviar un $u$ no can\'onico; reutilizar un $sid$ ya consumido; usar un $sid$
vencido; pedir una versi\'on inexistente; pedir un \code{record\_id}
inexistente o un \'indice fuera de $[1,N]$. &
HTTP 400, 409, 404 y 422, en ese orden. El \'indice inv\'alido lo rechaza el
cliente sin enviar nada, porque el servidor nunca ve $i$. \\
Manipulaci\'on &
Cambiar un bit en un registro $C_j$ del cat\'alogo, en una encapsulaci\'on
$E_j$, en un nonce o en los datos asociados. &
\code{sha256\_root} detecta el cambio en el cat\'alogo y el tag AEAD rechaza
los dem\'as (\code{InvalidTag}). \\
Ataque de desfase &
Un proxy entre cliente y servidor sustituye $u$ por $\hat u = u \cdot v$
(Ej.~11.20); en una variante, adem\'as desplaza una posici\'on la lista de
encapsulaciones. &
El cliente no acepta ning\'un registro: abrir $E_i$ falla con
\code{InvalidTag} en ambas variantes. \\
\bottomrule
\end{tabularx}
\caption{Plan de pruebas funcionales, negativas y de manipulaci\'on.}
\label{tab:plan-pruebas}
\end{table}

\subsection{Metodolog\'ia de experimentaci\'on}
Para medir el costo de ocultar la selecci\'on se evaluar\'an cat\'alogos de
$N \in \{500, 1000, 2500, 5000\}$ registros de 256\,B, comparando el modo OT
con la l\'inea base HTTP en la misma m\'aquina: una estaci\'on x86\_64
(Intel Core i7 o AMD Ryzen con AES-NI, 16\,GB de RAM) con Python~3.11. Los
tiempos se toman con \code{time.perf\_counter\_ns()}, con 5 repeticiones por
configuraci\'on para reportar media y desviaci\'on est\'andar.

Un script (\code{bench.py}) registrar\'a en CSV: (1) latencia de extremo a
extremo; (2) tiempo de c\'omputo en cliente y servidor; (3) operaciones de
grupo (cliente: 3 exponenciaciones y 1 multiplicaci\'on; servidor: $N+1$ y
$N$); y (4) bytes de la petici\'on, la respuesta y la descarga inicial. El
Cuadro~\ref{tab:plantilla-experimentos} adelanta los tama\~nos esperados
seg\'un los formatos de la Secci\'on~\ref{sec:formatos}; el Entregable~2
los contrastar\'a con lo medido.

\begin{table}[ht]
\centering
\small
\begin{tabularx}{\textwidth}{@{}r X X X X X@{}}
\toprule
$N$ & Descarga inicial & Petici\'on OT & Respuesta OT & Respuesta normal & Latencia (ms) \\
\midrule
500  & $\approx 0.22$\,MB & 96\,B & $\approx 55$\,KB  & $\approx 0.3$\,KB & \emph{A medir en E2} \\
1000 & $\approx 0.44$\,MB & 96\,B & $\approx 111$\,KB & $\approx 0.3$\,KB & \emph{A medir en E2} \\
2500 & $\approx 1.10$\,MB & 96\,B & $\approx 279$\,KB & $\approx 0.3$\,KB & \emph{A medir en E2} \\
5000 & $\approx 2.19$\,MB & 96\,B & $\approx 559$\,KB & $\approx 0.3$\,KB & \emph{A medir en E2} \\
\bottomrule
\end{tabularx}
\caption{Tama\~nos estimados de los cuerpos JSON (sin cabeceras HTTP) y plantilla para las latencias OT y normal del Entregable~2.}
\label{tab:plantilla-experimentos}
\end{table}

\section{Glosario de s\'imbolos}
\label{sec:glosario}

Para facilitar la lectura y unificar la notaci\'on utilizada a lo largo del
documento, en el Cuadro~\ref{tab:glosario-simbolos} se resume la nomenclatura
formal de par\'ametros y variables.

\begin{table}[!ht]
\centering
\small
\begin{tabularx}{\textwidth}{@{}l X@{}}
\toprule
S\'imbolo & Descripci\'on formal \\
\midrule
$G, g, q$ & Grupo algebraico c\'iclico (\texttt{ristretto255}), generador can\'onico $g$ y orden primo $q = 2^{252}+\dots$ \\
$N, i, j$ & Cantidad total de registros ($500 \dots 5000$); \'indice solicitado $i$; e \'indice de iteraci\'on $j$. \\
$K_j, C_j$ & Clave sim\'etrica AES-256 del registro $j$ y texto cifrado de 256\,B m\'as tag ($272$\,B en total). \\
$\beta, v$ & Escalar secreto ef\'imero del emisor ($\beta \in \mathbb{Z}_q$) y par\'ametro p\'ublico de sesi\'on ($v = g^\beta \in G$). \\
$\alpha, u$ & Escalar secreto del receptor ($\alpha \in \mathbb{Z}_q$) y elemento de consulta enmascarado ($u = g^\alpha v^{-i} \in G$). \\
$sid$ & Identificador fresco de sesi\'on de 128 bits (16\,bytes, representado en hexadecimal). \\
$\mathsf{ctx}$ & Cadena de contexto global inmutable con prefijos de longitud: cat\'alogo, versi\'on y grupo. \\
$AAD_{\mathrm{cat}}, AAD_{\mathrm{ot}}$ & Datos asociados: $\mathsf{ctx}\Vert j$ para registros y $\mathsf{ctx}\Vert sid\Vert j$ para encapsulaciones. \\
$u_j, Z_j, k_j, E_j$ & Clave p\'ublica impl\'icita ($u\cdot v^j$), secreto DH ($u_j^\beta$), clave derivada HKDF y clave envuelta. \\
$\mathrm{HKDF}, \mathrm{AEAD}$ & Funci\'on de derivaci\'on HKDF-SHA256 (RFC~5869) y cifrador autenticado AES-256-GCM (SP~800-38D). \\
\bottomrule
\end{tabularx}
\caption{Glosario formal de s\'imbolos y variables del protocolo.}
\label{tab:glosario-simbolos}
\end{table}


\bibliographystyle{plainnat}
\bibliography{refs}
\end{document}
